% SPDX-License-Identifier: LPPL-1.3c
% Copyright (C) 2026 Christophe Jorssen
% Author and Current Maintainer: Christophe Jorssen <christophe.jorssen@gmail.com>
% LPPL maintenance status: maintained
% This file is part of luafemm. See LICENSE and LICENSES.md for its terms.
\section{Measurement paths and the PGFPlots library}
\label{sec:profiles}
Load |\usepgfplotslibrary{femm}| after PGFPlots for the standard library
entry point. It loads the generic TikZ library if necessary. The profile
capture and command are also available through |femm| for compatibility;
only plotting requires PGFPlots.

\begin{key}{/tikz/femm/profile=\marg{options}}
Register an evaluated measurement path, without adding a physical region.
Use |\draw| to display the path or |\path| to keep it invisible. A single
connected open or closed path is required; arcs, circles, Bézier curves,
rounded corners, named coordinates and local transformations are accepted.
Do not combine |femm/region| and |femm/profile| on the same path.
Capture itself does not mesh or solve the model.
It may precede or follow physical declarations; only the first sampling
or plotting request requires those declarations to be complete. See
Section~\ref{sec:workflow}. Each profile remains one connected curve;
compound material paths require separate profiles for their separate contours.
\end{key}

The profile subkeys belong to |/femm/profile|:
\begin{key}{/femm/profile/name=\meta{identifier}}
Required name, with letters, digits, hyphens or underscores. The registry
is document-wide; reusing a name replaces its old geometry and cached
values. A profile keeps its own model even after another picture starts.
\end{key}
\begin{key}{/femm/profile/samples=\meta{integer} (initially 201)}
Number of equally spaced arc-length samples, including both endpoints;
between 2 and 100000. Increasing this number refines the graph's sampling,
not the finite-element solution. A narrow peak may be missed between samples.
\end{key}
\begin{key}{/femm/profile/curve tolerance=\meta{positive number} (initially .03)}
Flattening tolerance in model units. Independent of the material-path
tolerance. Samples follow the resulting polyline, so its arc length is
an approximation to the original smooth curve's length.
\end{key}
\begin{key}{/femm/profile/offset=\meta{number} (initially 0)}
Signed displacement towards the path's left normal, in model units.
Use a small nonzero offset to choose a side of a material interface.
The reported distance $s$ remains measured on the original path; the
sampled $x,y$ coordinates include the displacement. At sharp vertices,
offsets follow the chosen segment normal, not a rounded offset curve.
\end{key}
\begin{key}{/femm/profile/outside=\meta{error or nan} (initially error)}
PGF choice controlling samples outside the rectangular domain. |error|
stops with the offending point. |nan| supplies missing values, creating
gaps in the graph. It does not hide a mesher or solver failure.
\end{key}

\begin{command}{\femmplot\oarg{options}}
Inside a PGFPlots axis, emit a complete |\addplot+ ... coordinates ...;|
command. Do not append another semicolon. The first request solves the
profile's source model if necessary and caches all components. Further
component plots reuse those rows. Values pass in memory; no data file,
Python process or shell escape is involved. Options are parsed locally;
|\addplot+| executes in the enclosing axis so that legends and cycle lists
advance normally.
\end{command}

Every call resets the following plot options, in family |/femm/plot|:
\begin{key}{/femm/plot/profile=\meta{identifier}}
Required name of a previously registered profile.
\end{key}
\begin{key}{/femm/plot/component=\meta{choice} (initially B)}
PGF choice: |B|, |Bx|, |By|, |Bt|, |Bn|, |H|, |Hx|, |Hy|, |Ht|, |Hn| or |Az|.
Unsubscripted B and H mean the magnitudes. B components are in T, H
components in A/m and Az in T\,m. Cartesian components refer to the model
frame. Tangential and normal components follow the path direction.
\end{key}
\begin{key}{/femm/plot/abscissa=\meta{choice} (initially s)}
|s| is distance from the first point in model units; |t| is $s/L$ from
zero to one; |x| and |y| are sampled coordinates. In the TikZ interface,
distances are in millimetres. The parameter |t| is \emph{not} the polynomial
parameter of a Bézier curve. Nonmonotone x/y paths are plotted in traversal
order, not sorted by the chosen abscissa.
\end{key}
\begin{key}{/femm/plot/plot options=\marg{PGFPlots options}}
Ordinary options appended to |\addplot+|, empty by default. For example
|{red,thick,dashed,no markers}|. All axis labels, bounds, grids, legends and
cycle lists remain PGFPlots options. The emitted plot also sets
|unbounded coords=jump| unless explicitly overridden here.
\end{key}

\clearpage
\subsection{A curve and its components}
This example includes both the source picture and the graph and is
executed during the manual build.
\begin{codeexample}[width=6cm]
\begin{tikzpicture}[femm/problem={
  xmin=-15,xmax=15,ymin=-15,ymax=15,
  mesh size=2}]
  \filldraw[femm/region={
    material=cast-alnico-5-lng37,
    magnetization angle=90},
    fill=red!20]
    (-3,-4) rectangle (3,4);
  \draw[femm/profile={name=demo,
    samples=51},blue,dashed]
    (-8,6) .. controls (0,12) .. (8,6);
\end{tikzpicture}
\par\medskip
\begin{tikzpicture}
\begin{axis}[width=5.8cm,height=4.5cm,
  xlabel={$s$ (mm)},ylabel={Field (T)},
  no markers]
  \femmplot[profile=demo,component=B,
    plot options={black,thick}]
  \addlegendentry{$|\mathbf B|$}
  \femmplot[profile=demo,component=By,
    plot options={red}]
  \addlegendentry{$B_y$}
\end{axis}
\end{tikzpicture}
\end{codeexample}

\subsection{Complete example: straight and curved scans}
This example combines catalogue iron, alnico and passive copper. The first
scan traverses both air gaps and plots Cartesian components. The second
follows a cubic curve and plots tangent and normal components. Both graphs
retain the finite-element jumps rather than applying graphical smoothing.
\luafemmexample{materials-profiles}

\subsection{Tangents, interfaces and plotting conventions}
For unit tangent $\hat{\mathbf t}$, the left normal is
$\hat{\mathbf n}=(-\hat t_y,\hat t_x)$.
The signed components are $B_t=\mathbf B\cdot\hat{\mathbf t}$ and
$B_n=\mathbf B\cdot\hat{\mathbf n}$, and likewise for H. Reversing the
path changes these signs at a common point, while Bx and By are unchanged.
At a polygon vertex, the following segment supplies the tangent; the
final point uses the last segment. A closed path repeats its first
position as its final sample, but endpoint tangents can differ at a corner.

The P1 potential is continuous. Its gradient, and therefore B and H,
is constant in each triangle and may jump between elements or materials.
No smoothing is applied. PGFPlots normally joins consecutive samples by
straight segments. |const plot| changes the graphic convention but does
not locate all triangle-edge intersections. At an exact interface the
first containing triangle supplies the value; an offset or a path wholly
inside the desired medium makes the sampling side explicit.
